\documentclass[11pt]{article}
\usepackage[margin=0.9in]{geometry}
\usepackage{amsmath,amssymb,amsthm,booktabs,mathrsfs}
\usepackage[T1]{fontenc}
\usepackage{lmodern,microtype}
\usepackage[colorlinks=true,linkcolor=blue,urlcolor=blue]{hyperref}
\newtheorem{theorem}{Theorem}[section]
\newtheorem{lemma}[theorem]{Lemma}
\newtheorem{proposition}[theorem]{Proposition}
\newtheorem{corollary}[theorem]{Corollary}
\theoremstyle{remark}
\newtheorem{remark}[theorem]{Remark}
\newcommand{\norm}[1]{\left\|#1\right\|}
\title{Toward a four-dimensional Yang--Mills construction\\
\large Notebook 27: an exact composite blocking step}
\author{Research working document}
\date{17 October 2026}
\begin{document}
\maketitle
\begin{abstract}
Fix one primitive ordinary averaging ratio $\ell$ and take a finite
macro ratio $R=\ell^s$. We construct the local chart and Gaussian
of the actual composite constraint $M_\ell^s$, using a separately
chosen direct-$R$ internal contour gauge. A selected interface bond
has exactly one nonlinear consumer at every intermediate level.
Its flat pivot coefficient is $R^{1-d}$, and elimination preserves
all retained coordinates. A finite rectangular-patch argument gives
coercivity of curvature, contour gauge and the composite constraint;
the preceding classical action bounds then yield a retained Gaussian
gap and exponentially decaying covariance uniformly in the old
refinement. We derive the compact gauge-orbit normalization and the
local Haar, delta, pivot, coupling and Gaussian factors of the
one-shot constrained integral. Two minimum comparisons identify its
physical centered branch. Its pre-gauge representative supplies a
common-ball continuation of an inherited gauge-invariant insertion
without changing integration variables or its cutoff.
The construction is a finite small-field step. Uniform complex
nonvanishing, supported inherited interaction bounds, amplitude
control and the remaining sectors of the existence proof require
their separate estimates.
\end{abstract}

The purpose of this construction is to make the factor $R^{-1}$ of
Notebook 25 available in an actual blocking integral. The primitive
ratio remains fixed while $R=\ell^s$ is chosen large enough for the
desired representative bound. The constraint is defined by its
composition; no identity with a separately defined direct-$R$
ordinary average is needed. Likewise, this one-shot integration is
not identified with $s$ already truncated primitive integrations.
The gauge choice and every factor of the new integral are specified
below. All old-depth uniformity refers to $M=\ell^k$, with $\ell,R$
fixed and compatible periodic volume.


\section{Geometry of the actual composite ordinary constraint}
\label{cmg-section}

Fix an odd primitive ratio $\ell\ge3$, a positive integer $s$, and
$R=\ell^s$. Current spacing is one. An earlier classical refinement
$M=\ell^k$ does not enter any map in this section. Constants are uniform
in $M$ and compatible periodic volume once $\ell,s$ are fixed; no
uniformity in $R\to\infty$ is asserted. Compatible source periods
$2\ell^n$, $n\ge s$, suffice. Local paths are evaluated in their specified
lifts, so the two opposite interfaces of a two-block periodic direction
remain distinct. Translate the common block center to zero and identify
the level-$t$ lattice with $\ell^t\mathbb Z^d$.

The constraint constructed here is the actual composition of $s$
ordinary primitive averages. It is not identified with a direct-$R$
ordinary average or with $s$ separately truncated quantum integrals.
The direct-$R$ contour variable used below is a separate gauge choice.

\subsection{The source map and its exact covariance}
\label{cmg-source-section}

For a centered primitive block and $x\in B_\ell(y)$, let
$\mathcal G_\ell(y,x)$ contain the shortest coordinate-order paths
from $y$ to $x$. Every such path is internal to this block. Use the
selected source group mean $\mathfrak M$ to define
\begin{equation}
 P_\ell(V;y,x)=
 \mathfrak M\bigl(\{V(\Gamma):\Gamma\in\mathcal G_\ell(y,x)\}\bigr).
 \label{cmg-contour}
\end{equation}
For $c=(y,y+\ell e_\mu)$, put $x'=x+\ell e_\mu$ and let
$R_c(V)$ be the straight center transporter. The ordinary map is
\begin{align}
 U_x(V)&=P_\ell(V;y,x)V([x,x'])P_\ell(V;y+\ell e_\mu,x')^{-1},
 \notag\\
 \mathcal M_\ell(V)(c)
 &=\exp\left\{\ell^{-d}\sum_{x\in B_\ell(y)}
          \log[U_x(V)R_c(V)^{-1}]\right\}R_c(V).
 \label{cmg-ordinary}
\end{align}
All logarithms use one common small branch. These are the ordinary
average and contour definitions in \cite{balaban109cm}, pp.~252--254,
(0.3), (0.5)--(0.12).

The source axioms give local analyticity, the arithmetic differential
at equal identity inputs, and endpoint covariance
\begin{equation}
 \mathfrak M(\{gU_jh^{-1}\})=
 g\mathfrak M(\{U_j\})h^{-1}.
 \label{cmg-mean-covariance}
\end{equation}
For completeness, (0.10) of that paper supplies one admissible
realization as the unique small solution $P$ of
$\sum_j\log(U_jP^{-1})=0$. The derivative with respect to $\log P$ at
identity is minus the number of inputs times the identity, so the
analytic implicit-function theorem gives uniqueness, analyticity and
the arithmetic differential. Conjugating the defining logarithms
proves \eqref{cmg-mean-covariance} for this realization. The paper
explicitly permits any mean satisfying its axioms; the proof here
uses those properties for the selected mean and does not replace it
by a different one.

With $V^u(b)=u(b_-)V(b)u(b_+)^{-1}$, it follows that
\begin{align}
 P_\ell(V^u;y,x)&=u(y)P_\ell(V;y,x)u(x)^{-1},\notag\\
 \mathcal M_\ell(V^u)&=
       \mathcal M_\ell(V)^{\,u|_{\text{new centers}}}.
 \label{cmg-primitive-covariance}
\end{align}
Indeed every relative transporter inside the logarithm in
\eqref{cmg-ordinary} is conjugated at its source center. The identities
hold on the common analytic domains and extend to complex fields there.

\subsection{One primitive pivot has one nonlinear consumer}
\label{cmg-primitive-section}

Select the unit interface bond
\begin{equation}
 p_\ell(c)=
 [y+(\ell-1)e_\mu/2,\ y+(\ell+1)e_\mu/2].
 \label{cmg-primitive-pivot}
\end{equation}
No internal contour reads this bond. A parallel string containing it
must have direction $\mu$, the same transverse block center, and
the same pair of longitudinal endpoint blocks. Hence it belongs to
the output equation $c$. The reference center string has the same
property. This is a statement about the actual products and group
means before differentiation:
\begin{equation}
 \mathcal M_\ell(V)(c')\text{ is independent of }V(p_\ell(c))
 \quad(c'\ne c).
 \label{cmg-primitive-incidence}
\end{equation}
The local lifted strings distinguish opposite periodic interfaces.

Let $C_\ell$ be the angle differential at identity, and let
$t_y(B;x)$ be the arithmetic mean of the oriented sums of $B$ on
the root-to-$x$ contours. Linearizing \eqref{cmg-ordinary} gives
\begin{equation}
 C_\ell B(c)=\ell^{-d}\sum_{x\in B_\ell(y)}
 \left[t_y(B;x)+B([x,x+\ell e_\mu])
             -t_{y+\ell e_\mu}(B;x+\ell e_\mu)\right].
 \label{cmg-primitive-linear}
\end{equation}
The reference-string derivative cancels between the logarithms and
the final factor. The root terms do not read the pivot. Exactly $\ell$
starting sites have its transverse coordinates, and each corresponding
positive string contains it once. Thus its full column is
\begin{equation}
 (C_\ell)_{c',p_\ell(c)}
       =\ell^{1-d}\delta_{c,c'}I_{\mathfrak g}.
 \label{cmg-primitive-coefficient}
\end{equation}
This is an angle derivative, without division by the coarse length.

\subsection{Exact nonlinear incidence through the composition}
\label{cmg-composite-section}

Define
\begin{equation}
 \mathcal M_R^{\mathrm{comp}}=\mathcal M_\ell^s.
 \label{cmg-composite}
\end{equation}
For a final bond $c=(y,y+Re_\mu)$ choose
\begin{equation}
 p_R(c)=[y+(R-1)e_\mu/2,\ y+(R+1)e_\mu/2].
 \label{cmg-macro-pivot}
\end{equation}
Writing $h_t=\ell^t$, set
\begin{equation}
 b_t(c)=[y+(R-h_t)e_\mu/2,\ y+(R+h_t)e_\mu/2],
 \qquad 0\le t\le s.
 \label{cmg-pivot-chain}
\end{equation}
Its endpoints are level-$t$ sites because $R/h_t$ is odd.
Moreover $b_0=p_R(c)$, $b_s=c$, and $b_t$ is exactly the primitive
middle-interface pivot of $b_{t+1}$.

\begin{proposition}
\label{cmg-incidence-proposition}
Changing the input group coordinate at $p_R(c)$ changes at level $t$
at most $b_t(c)$. Consequently every final constraint reads its own
selected pivot and no other selected pivot. Each final output reads
only unit bonds whose endpoints lie in its two endpoint $R$-blocks,
using the prescribed local lift. If $C_R$ denotes the composite angle
differential at identity, then
\begin{equation}
 (C_R)_{c,p_R(c')}=R^{1-d}\delta_{c,c'}I_{\mathfrak g}.
 \label{cmg-composite-coefficient}
\end{equation}
\end{proposition}

\begin{proof}
The incidence assertion is induction along
\eqref{cmg-pivot-chain}, applying
\eqref{cmg-primitive-incidence} at spacing $h_t$.
It is functional independence, so is valid for nonlinear variations,
not only for the derivative. Distinct final bonds have distinct
intermediate selected bonds.

For locality, a level-$t$ output reads level-$(t-1)$ bonds in its two
endpoint primitive blocks, including their connecting strings.
The descendants centered at those intermediate sites tile the two
centered blocks of side $\ell^t$. Iterating gives the stated unit-bond
dependency set. On short compatible tori this is read in the local
cover, not inferred merely from periodic set containment.

Finally the chain rule along the single affected output multiplies
$s$ copies of $\ell^{1-d}I$. There is no intermediate row sum.
The product is $R^{1-d}I$, proving
\eqref{cmg-composite-coefficient}.
\end{proof}

Composition of \eqref{cmg-primitive-covariance} gives
\begin{equation}
 \mathcal M_R^{\mathrm{comp}}(V^u)=
 \mathcal M_R^{\mathrm{comp}}(V)^{\,u|_{R\text{-centers}}}.
 \label{cmg-composite-covariance}
\end{equation}
With $D\phi(b)=\phi(b_+)-\phi(b_-)$, differentiate using
$u=e^{-it\phi}$ to obtain
\begin{equation}
 C_RD\phi(c)=\phi(c_+)-\phi(c_-).
 \label{cmg-gradient}
\end{equation}
The right side is an endpoint difference, or $R$ times the physical
coarse gradient.

\subsection{Local analytic elimination and unchanged retained coordinates}
\label{cmg-chart-section}

For a small background $V=e^{ia}$ define
\begin{align}
 F_{R,V}(B;c)&=\frac1i\log\left[
 \mathcal M_R^{\mathrm{comp}}(e^{iB}V)(c)
 \mathcal M_R^{\mathrm{comp}}(V)(c)^{-1}\right],\notag\\
 A_{R,V}&=D_BF_{R,V}(0).
 \label{cmg-relative-map}
\end{align}
Finite composition and the two-block dependency give a positive
product domain in $a,B$, independent of volume and old $M$ at fixed
$R$. Put
\begin{equation}
 \pi_c(V)=R^{d-1}(A_{R,V})_{c,p_R(c)}.
 \label{cmg-pi}
\end{equation}
It is local analytic, $\pi_c(1)=I$, and on a common smaller domain
$\lVert\pi_c-I\rVert\le1/2$, $\lVert\pi_c^{-1}\rVert\le2$.
The right inverse inserts only at selected pivots:
\begin{equation}
 (h_{R,V}w)(p_R(c))=R^{d-1}\pi_c(V)^{-1}w(c),\qquad
 (h_{R,V}w)(b)=0\ (b\text{ nonpivot}),\qquad A_{R,V}h_{R,V}=I.
 \label{cmg-right-inverse}
\end{equation}
In particular its maximum row norm is at most $2R^{d-1}$.
If $Q_{R,V}=A_{R,V}/R$ is the physical coarse-potential derivative,
and its kernel is written as
$QB(c)=\sum_bR^{-d}Q(c,b)B(b)$, that kernel's pivot coefficient
is exactly $\pi_c(V)$. Thus the factors in
\eqref{cmg-right-inverse} retain the source physical normalization.

Let $E_f$ insert nonpivot values and zero at pivots, and abbreviate
$A=A_{R,V}$, $h=h_{R,V}$. Define
\begin{equation}
 K_{R,V}=E_f-hAE_f,\qquad AK_{R,V}=0.
 \label{cmg-insertion}
\end{equation}
Restriction of $K_{R,V}z$ to nonpivots is $z$. Write
$F_{R,V}(B)=AB+\mathcal C_V(B)$, with
$\mathcal C_V(0)=D\mathcal C_V(0)=0$, and solve
\begin{equation}
 \widetilde{\mathcal D}_V(B)
       =\mathcal C_V(B-h\widetilde{\mathcal D}_V(B)).
 \label{cmg-implicit-correction}
\end{equation}
For output $c$ the right side reads only its own unknown
$\widetilde{\mathcal D}_V(B;c)$: other unknowns insert at pivots that
the full nonlinear $c$ equation does not read. The implicit derivative
at zero is identity. A finite-dimensional analytic implicit theorem,
or contraction after imposing $\lVert D\mathcal C_V h\rVert<1/2$,
therefore gives a unique local analytic solution. There are finitely
many oriented neighborhood types, so one radius works in all volumes.
Iteration does not spread its spatial dependency.

The exact retained chart is
\begin{align}
 \Phi_{R,V}(z)
   &=K_{R,V}z-h\widetilde{\mathcal D}_V(K_{R,V}z),\notag\\
 F_{R,V}(\Phi_{R,V}(z))&=0,\qquad
 \Phi_{R,V}(z)(b)=z(b)\quad(b\text{ nonpivot}).
 \label{cmg-retained-chart}
\end{align}
It is the unique small solution with its prescribed retained values.
Direct differentiation gives
\begin{align}
 D\Phi_{R,V}(0)&=K_{R,V},\notag\\
 D^2\Phi_{R,V}(0)[v,w]
   &=-hD^2F_{R,V}(0)[K_{R,V}v,K_{R,V}w].
 \label{cmg-chart-derivatives}
\end{align}
The affine equation $F_{R,V}(B)=w$ is equally local:
\begin{equation}
 B=K_{R,V}z+hw-h\mathcal C_V(B).
 \label{cmg-affine-chart}
\end{equation}
It has a common positive small product domain in $a,w,z$ and still
retains $z$ exactly. Local finite dimension, Cauchy estimates and
bounded neighborhood incidence give finite weighted first-derivative
row and column bounds and finite contracted second-derivative bounds.
Their constants may depend on $R$, but not on $M$ or volume.

After coupling insertion the relative field is
$\Phi_{R,V}(gz)$, whose retained values equal $gz$.
It is not $g\Phi_{R,V}(z)$. Additional pivot, contour or reconstructed
field indicators remain at their exact nonlinear arguments.

\subsection{An independent direct-$R$ contour gauge}
\label{cmg-gauge-section}

Use the same selected source mean on root-to-site coordinate-order
paths inside one centered $R$-block:
\begin{equation}
 P_R(V;y,x)=
 \mathfrak M\bigl(\{V(\Gamma):\Gamma\in\mathcal G_R(y,x)\}\bigr).
 \label{cmg-direct-contour}
\end{equation}
This does not redefine $\mathcal M_R^{\mathrm{comp}}$. Endpoint
covariance gives
\begin{equation}
 P_R(V^u;y,x)=u(y)P_R(V;y,x)u(x)^{-1}.
 \label{cmg-direct-covariance}
\end{equation}
Every path is internal to its block; every selected pivot crosses a
block interface. Thus $P_R$ is exactly independent of all selected
pivot group variables at every allowed background.

Define
\begin{equation}
 T_{R,V}B(y,x)=
 D_B\left.\frac1i\log[
 P_R(e^{iB}V;y,x)P_R(V;y,x)^{-1}]\right|_{B=0}.
 \label{cmg-T-definition}
\end{equation}
The pivot support of $h_{R,V}$ then gives
\begin{equation}
 T_{R,V}h_{R,V}=0,\qquad
 T_{R,V}K_{R,V}=T_{R,V}E_f.
 \label{cmg-T-identities}
\end{equation}
More strongly,
\begin{equation}
 P_R(e^{i\Phi_{R,V}(z)}V;y,x)
       =P_R(e^{iE_fz}V;y,x).
 \label{cmg-gauge-retained}
\end{equation}
The contour penalty therefore reads exactly the unchanged retained
variables, although its coefficients still depend on $V$.

For normalized trace $\tau$ put
\begin{equation}
 \mathcal G_R(V)=\sum_y\sum_{\substack{x\in B_R(y)\\x\ne y}}
       \left[1-\tfrac12\tau(P_R(V;y,x)+P_R(V;y,x)^{-1})\right].
 \label{cmg-gauge-action}
\end{equation}
It is nonnegative on real unitary fields. At a background with every
$P_R=I$ its first derivative vanishes and its pulled-back Hessian,
in counting/trace coordinates, is
\begin{equation}
 D_z^2\mathcal G_R(e^{i\Phi_{R,V}(z)}V)\big|_{z=0}
       =K_{R,V}^{T}T_{R,V}^{T}T_{R,V}K_{R,V}.
 \label{cmg-gauge-hessian}
\end{equation}
The second derivative of the chart contributes nothing because it
is pivot supported and the gauge action is pivot independent. At
identity, \eqref{cmg-direct-covariance} also gives
\begin{equation}
 T_{R,1}D\phi(y,x)=\phi(x)-\phi(y).
 \label{cmg-T-gradient}
\end{equation}

\subsection{Clearing this gauge and the domain boundary}
\label{cmg-clearing-section}

Set $u_V(y)=I$ at all final centers and
$u_V(x)=P_R(V;y,x)$ elsewhere in its block. Then, on the common
domain,
\begin{equation}
 P_R(V^{u_V};y,x)=I,\qquad
 \mathcal M_R^{\mathrm{comp}}(V^{u_V})
       =\mathcal M_R^{\mathrm{comp}}(V).
 \label{cmg-clearing}
\end{equation}
This is a geometric gauge operation, not by itself a minimum or
stationarity argument.

There are at most $d!$ contour orderings and every path has length
at most $d(R-1)/2$. In a fixed matrix norm,
$\lVert\prod_b e^{ia_b}-I\rVert\le
\exp(\sum_b\lVert a_b\rVert)-1$. The selected local mean has a
fixed Lipschitz constant near equal identity inputs, with the number
of inputs bounded by $d!$. Consequently, when $R\lVert a\rVert_\infty$
is sufficiently small,
\begin{equation}
 \lVert\log u_V/i\rVert_\infty
       \le C R\lVert a\rVert_\infty.
 \label{cmg-gauge-size}
\end{equation}
The constant here depends on dimension, group and the selected mean,
not on $R$. For a pre-gauge representative
$a_p=O(w/R)$ this gives $u=O(w)$. It does not give a rooted link-angle
bound $O(w/R)$: endpoint gauge differences may have size $O(w)$.
Nor does \eqref{cmg-gauge-size} make the local elimination radius
uniform in $R$.

The composite constraint, its actual local chart, and the independent
contour gauge are thus established. Gaussian coercivity and locality,
the physical minimum comparisons, the exact compact gauge insertion
and all cutoff normalizations are separate modules. None of these
geometric identities supplies a repeated quantum RG induction.

\section{The retained Gaussian for the composite macrostep}
\label{cmq-section}

Fix an odd primitive ratio $\ell\geq3$, a finite macro ratio
$R=\ell^s$, $s\geq1$, and an arbitrary preceding classical depth
$M=\ell^k$, $k\geq0$. Current lattice spacing is one. All constants
in this section may depend on the fixed $\ell,s$, but are independent
of $M$ and compatible periodic volume. We use the actual composite
ordinary constraint $F_R=\mathcal M_\ell^s$ and an independently chosen
direct-$R$ internal contour gauge. No equality with a direct-$R$
ordinary average is used.

\subsection{Established classical inputs and the actual flat constraint}
\label{cmq-input-section}

Let $S_M(a)$ be the homogeneous ordinary classical minimum action
at current datum $e^{ia}$. Its flat Hessian $q_M$ has the previously
established variational description
\begin{equation}
 \langle B,q_MB\rangle
 =\inf_{Q_M^{\rm st} A=B}\|d_{1/M}A\|_{1/M}^2,
 \qquad
 \|d_1B\|_2^2\leq\langle B,q_MB\rangle
 \leq C_d\|B\|_2^2,
 \quad C_d=2^{d+6}(d-1).
 \label{cmq-classical-input}
\end{equation}
The case $M=1$ is included. In current counting norms, $q_M$ has a
uniform weighted Schur bound at some positive exponent $\kappa_q$.
On a common independent-background product ball, the earlier
classical construction also gives
\begin{equation}
 DS_M(0)=0,\qquad
 \|DS_M(a)\|_\infty\leq C\|a\|_\infty,
 \qquad
 \|D^2S_M(a)-q_M\|_\kappa\leq C_\kappa\|a\|_\infty
 \label{cmq-classical-nonlinear-input}
\end{equation}
at a sufficiently small positive exponent. These are inputs from
the homogeneous classical action; their constants do not depend
on the next chosen macro ratio.

Define the composite \emph{angle} differential
\begin{equation}
 C_RB=\left.\frac{d}{dt}\frac1i
            \log F_R(e^{itB})\right|_{t=0}.
 \label{cmq-constraint}
\end{equation}
It is not divided by $R$. The exact composite geometry supplies
the following facts. A row for the final bond $c=(y,y+Re_\mu)$
reads only unit bonds in the union of its two endpoint $R$-blocks.
Endpoint gauge covariance gives
\begin{equation}
 C_RD\phi(c)=\phi(y+Re_\mu)-\phi(y),
 \qquad D\phi(b)=\phi(b_+)-\phi(b_-).
 \label{cmq-gradient}
\end{equation}
Choose the unit centerline interface pivot
\begin{equation}
 p(c)=\big[y+(R-1)e_\mu/2,\ y+(R+1)e_\mu/2\big].
 \label{cmq-pivot}
\end{equation}
At level $t$ its unique affected intermediate bond has endpoints
$y+(R-\ell^t)e_\mu/2$ and $y+(R+\ell^t)e_\mu/2$.
Thus the primitive pivot coefficient $\ell^{1-d}$ multiplies
along one chain and yields
\begin{equation}
 (C_R)_{c,p(c')}=R^{1-d}\delta_{cc'}I.
 \label{cmq-pivot-coefficient}
\end{equation}
These statements hold in the periodic covering geometry with
the specified lifted paths and distinct opposite interfaces.

Let $E_f$ insert nonpivot coordinates and zero at pivots. Define
\begin{equation}
 (h_Rw)(p(c))=R^{d-1}w(c),\qquad
 K_R=(I-h_RC_R)E_f.
 \label{cmq-retained-map}
\end{equation}
Then $C_Rh_R=I$, $C_RK_R=0$, $K_R$ retains its input coordinates,
and $\operatorname{ran}K_R=\ker C_R$. Both maps have fixed
$R$-range. One conservative finite norm bound is
\begin{equation}
 r_\ell=1+\frac{2d(\ell-1)}\ell,\qquad
 c_R=\big[\ell r_\ell\sqrt{2d+2}\big]^s,\qquad
 \|C_R\|\leq c_R,\qquad
 \|K_R\|^2\leq k_R^2:=1+R^{2d-2}c_R^2.
 \label{cmq-retained-upper}
\end{equation}
Indeed the primitive counting norm is bounded by
$\ell r_\ell\sqrt{2d+2}$ and these angle spaces compose.
The retained and pivot parts of $K_Rz$ have disjoint coordinate
supports, which proves the last bound.

\subsection{The independent contour gauge and a local nullspace argument}
\label{cmq-patch-section}

Use the actual direct-$R$ internal averaged contours solely as
gauge fixing, with flat derivative
\begin{equation}
 (T_RB)(y,x)=\operatorname{av}_{\Gamma:y\to x}
              \sum_{b\in\Gamma}\operatorname{or}(b)B_b.
 \label{cmq-contour}
\end{equation}
Every contour stays in its own centered $R$-block. Consequently
\begin{equation}
 T_RD\phi(y,x)=\phi(x)-\phi(y),\qquad T_Rh_R=0.
 \label{cmq-contour-gradient}
\end{equation}
For normalized trace $\tau$, the absolute real gauge action
$G_R(V)=\sum_{y,x}[1-\operatorname{Re}\tau(P_R(V;y,x))]$
has zero constant and linear terms at $V=1$, and flat Hessian
$T_R^TT_R$. This follows from
$P_R(e^{itB})=I+itT_RB+O(t^2)$ and unitarity, or from its
local logarithm. In the complex domain use
$1-\tau(P+P^{-1})/2$. The elementary row/column estimate gives
\begin{equation}
 \|T_R\|^2\leq t_R^2:=d(R-1)(R^d-1).
 \label{cmq-contour-upper}
\end{equation}
The two factors bound contour length and the number of target
rows through an internal bond. This is the actual averaged
contour penalty, not an assumed tree-coordinate penalty.

Let $\Omega$ be the rectangular cube consisting of the $3^d$
centered $R$-blocks indexed by $\{-1,0,1\}^d$. Include all
positive unit bonds with both endpoints in $\Omega$ and all
unit plaquettes with vertices there. Let $T_{R,\Omega}$ contain
all internal contour rows in these blocks, and let
$C_{R,\Omega}$ contain the neighboring-center rows whose two
centers belong to this cube. The two-block locality places every
one of these rows inside the chosen bond set. Set
\begin{equation}
 \mathcal L_{R,\Omega}
 =d_\Omega^Td_\Omega+T_{R,\Omega}^TT_{R,\Omega}
                         +C_{R,\Omega}^TC_{R,\Omega},
 \qquad \sigma_R=\lambda_{\min}(\mathcal L_{R,\Omega}).
 \label{cmq-local-matrix}
\end{equation}

\begin{lemma}[Local positivity]
\label{cmq-local-positivity}
The finite scalar matrix in \eqref{cmq-local-matrix} has
$\sigma_R>0$.
\end{lemma}
\begin{proof}
If its quadratic form vanishes, $d_\Omega B=0$. The rectangular
grid is simply connected: its plaquette relations imply path
independence and hence $B=D\phi$ on every bond of $\Omega$.
Equation \eqref{cmq-contour-gradient} makes $\phi$ constant on
each $R$-block. Equation \eqref{cmq-gradient} makes the constants
on neighboring blocks equal. Its local use is justified by
extending $\phi$ arbitrarily outside $\Omega$ and then using
two-block locality of the row. The coarse block graph is
connected, so $\phi$ is constant on all of $\Omega$ and $B=0$.
The matrix therefore has trivial kernel and is positive definite.
\end{proof}

This explicitly excludes the gauge and harmonic obstructions.
Mere finiteness of a local matrix would not have supplied the
required positivity.

Translate the patch over all $R$-block centers. Every bond occurs
at least once, whereas every plaquette, contour row and
constraint row occurs at most $3^d$ times. Thus
\begin{equation}
 \sigma_R\|B\|_2^2
 \leq3^d\big(\|d_1B\|_2^2+\|T_RB\|_2^2+\|C_RB\|_2^2\big).
 \label{cmq-global-patch}
\end{equation}
On a compatible torus, pull the periodic field back to the
ordinary covering cube and retain repeated occurrences. Summing
the same local inequalities has exactly the same multiplicity
bound; an embedded nonwrapping patch is not required. The
distinct-interface assumptions of the composite chart remain
in force.

Combining \eqref{cmq-classical-input} and
\eqref{cmq-global-patch}, for $B\in\ker C_R$ we obtain
\begin{equation}
 \langle B,q_MB\rangle+\|T_RB\|_2^2
 \geq\gamma_R\|B\|_2^2,\qquad
 \gamma_R=\sigma_R/3^d>0.
 \label{cmq-field-gap}
\end{equation}
The upper bound on arbitrary $B$ is
$(C_d+t_R^2)\|B\|_2^2$. None of these constants depends on $M$.

\subsection{Precision, covariance and weighted inverse}
\label{cmq-inverse-section}

At flat background both scalar gradients vanish. Therefore the
second derivative of the nonlinear constraint chart contributes
no gradient-times-chart-curvature term. The exact retained flat
precision is
\begin{equation}
 A_{M,R,0}=K_R^T(q_M+T_R^TT_R)K_R.
 \label{cmq-flat-precision}
\end{equation}
Since retained restriction of $K_Rz$ is $z$,
\begin{equation}
 \gamma_R I\leq A_{M,R,0}\leq B_R I,\qquad
 B_R=(C_d+t_R^2)k_R^2,\qquad
 B_R^{-1}I\leq A_{M,R,0}^{-1}\leq\gamma_R^{-1}I.
 \label{cmq-sandwich}
\end{equation}
There are $d(R^d-1)\dim\mathfrak g$ retained real coordinates
per $R$-block. Their normalized Gaussian has density
\begin{equation}
 (2\pi)^{-n/2}\det(A_{M,R,0})^{1/2}
        \exp[-\tfrac12 z^TA_{M,R,0}z]\,dz.
 \label{cmq-density}
\end{equation}
Its covariance is the counting inverse $A_{M,R,0}^{-1}$.
There is no original-fine volume factor $M^d$ in these current
coordinates. Haar and nonlinear Jacobian terms remain in the
interaction.

For a block matrix, write $\|A\|_{{\rm Sch},\alpha}$ for the
maximum of its exponentially weighted row and column sums of
block operator norms. Regrouping the current $q_M$ bound into
$R$-blocks gives, for $\alpha/R\leq\kappa_q$,
\begin{equation}
 \|q_M\|_{{\rm Sch},\alpha;R}
 \leq R^d e^{\alpha d} K_q(\alpha/R).
 \label{cmq-regroup}
\end{equation}
The fixed-range maps $K_R,T_R$ and weighted matrix
multiplication then give some $\alpha_R>0$ and $A_R^*<\infty$
such that
\begin{equation}
 \|A_{M,R,0}\|_{{\rm Sch},\alpha_R}\leq A_R^*
 \label{cmq-precision-schur}
\end{equation}
uniformly in $M$ and compatible volume.

Here is an explicit inverse argument. Fix an $R$-block $j$ and
put $W_{ii}=e^{t d(i,j)}I$. Periodic distance is used on a torus.
For $0<t\leq\alpha_R/2$, both the row and column Schur bounds
of $WAW^{-1}-A$ are at most
\begin{equation}
 \sup_i\sum_k\|A_{ik}\|(e^{t d(i,k)}-1)
 \leq\frac{2tA_R^*}{e\alpha_R}.
 \label{cmq-conjugation}
\end{equation}
Indeed $e^{tx}-1\leq txe^{tx}$ and
$\sup_{x\geq0}xe^{-\alpha_Rx/2}=2/(e\alpha_R)$.
Choose
\begin{equation}
 t_R^{\rm loc}
 =\min\left\{\frac{\alpha_R}{2},
                  \frac{e\alpha_R\gamma_R}{4A_R^*}\right\}.
 \label{cmq-rate}
\end{equation}
The perturbation has Hilbert norm at most $\gamma_R/2$.
The resolvent identity and \eqref{cmq-sandwich} give
$\|(WAW^{-1})^{-1}\|\leq2/\gamma_R$, whence
\begin{equation}
 \|(A_{M,R,0}^{-1})_{ij}\|
 \leq\frac2{\gamma_R}e^{-t_R^{\rm loc}d(i,j)}.
 \label{cmq-entry-decay}
\end{equation}
On the infinite lattice, truncated distance weights give the
same estimates and may be removed in each matrix entry. At
$\kappa_R=t_R^{\rm loc}/2$ this implies
\begin{equation}
 \|A_{M,R,0}^{-1}\|_{{\rm Sch},\kappa_R}
 \leq\frac2{\gamma_R}
       \sum_{z\in\mathbb Z^d}e^{-t_R^{\rm loc}|z|_1/2}.
 \label{cmq-covariance-schur}
\end{equation}
Periodic sums are bounded by this same lattice sum by shortest
representatives. The proof uses both the real gap and the
already-established decay of $q_M$; a spectral sandwich alone
would not imply locality.

\subsection{The independent-background extension}
\label{cmq-background-section}

Let $\Phi_a(z)$ be the actual finite-range analytic pivot chart
for the composite relative constraint at $V=e^{ia}$, retaining
$z$ exactly. Define
\begin{equation}
 a_V(z)=\frac1i\log[e^{i\Phi_a(z)}e^{ia}],\qquad
 A_{M,R}(a)
 =\left.D_z^2\{S_M(a_V(z))+G_R(e^{i\Phi_a(z)}e^{ia})\}
                                                     \right|_{z=0}.
 \label{cmq-background-precision}
\end{equation}
The action contribution is exactly
\begin{equation}
 (Da_V)^T D^2S_M(a)Da_V+\langle DS_M(a),D^2a_V\rangle.
 \label{cmq-chain-rule}
\end{equation}
Thus both constraint-straightening curvature and ordinary-log
curvature are retained. The gauge contribution is its full
absolute Hessian, including its residual away from $P_R=I$.

The local chart derivatives have fixed-$R$ row/column bounds on
a smaller product ball. Contracting their local second
derivatives against the supremum bound for $DS_M(a)$, and using
\eqref{cmq-classical-nonlinear-input}, proves at a smaller fixed
exponent
\begin{equation}
 \|A_{M,R}(a)-A_{M,R,0}\|_{{\rm Sch},\kappa_R}
 \leq C_{R,\kappa_R}\|a\|_\infty.
 \label{cmq-background-bound}
\end{equation}
Every contraction is a finite local incidence sum; no extensive
gradient norm or old-$M$ coordinate count occurs. A sufficiently
small background ball gives complex accretivity and a
holomorphic weighted inverse by its Neumann series. For example,
one may require both
$C_{R,\kappa_R}\|a\|_\infty<\gamma_R/2$ and
$\|A_{M,R,0}^{-1}\|_{{\rm Sch},\kappa_R}
 C_{R,\kappa_R}\|a\|_\infty<1/2$.

These are uniform-old-depth Gaussian statements for the chosen
composite constraint and independent contour gauge. Physical
centering requires the separate stationary-branch comparison;
the exact Haar, delta, gauge-insertion and cutoff normalization
requires the separate integral argument. Strict background
carrier coefficients require the classical support composition,
not merely \eqref{cmq-background-bound}. This section proves
neither a quantum induction nor a uniform estimate as $R$ tends
to infinity.

% No preamble. Composite macrostep integral, unique prefix cmi-.
\section{The exact composite integral and its pre-gauge continuation}
\label{cmi-section}

Fix the primitive odd ratio $\ell$ and a macro ratio $R=\ell^s$.
The preceding refinement is $M=\ell^k$. All statements concern
finite compatible periodic regulators. Constants may depend on the
fixed $R,\ell$, but are uniform in $M$ and volume. The constraint is
the actual composite
\[
 \mathcal F_R=\mathcal M_\ell^s.
\]
It is not replaced by a direct-$R$ ordinary average. The direct-$R$
internal contour variable $P_R$ is used independently for gauge
fixing. Nor is the one-shot integral below identified with $s$
separately gauge-fixed and truncated quantum integrations.

\subsection{A gauge insertion with its exact normalization}
\label{cmi-gauge-section}

Use normalized Haar measure on the compact matrix group $G$, and
the convention $V^u(b)=u(b_-)V(b)u(b_+)^{-1}$. The previously
proved local geometry gives
\begin{align}
 \mathcal F_R(V^u)
   &=\mathcal F_R(V)^{u|_{\text{final centers}}},\notag\\
 P_R(V^u;y,x)&=u(y)P_R(V;y,x)u(x)^{-1}.                 \label{cmi-covariance}
\end{align}
Every contour defining $P_R$ stays inside its $R$-block.
Put
\begin{equation}\label{cmi-penalty}
 v(h)=1-\operatorname{Re}\tau h,\qquad
 \mathcal G_R(V)=\sum_y\sum_{\substack{x\in B_R(y)\\x\ne y}}
                       v(P_R(V;y,x)).
\end{equation}
The analytic continuation of $v(h)$ is
$1-\tfrac12\tau(h+h^{-1})$. This is the absolute contour
penalty. It is not replaced by a relative penalty away from
a gauge-zero background.

A local logarithmic formula does not by itself allow integration
over every compact gauge transformation. One may work on a
gauge-invariant domain where the contour formula is defined on
complete orbits, or make the following explicit measurable
completion. For a path tuple $(U_j)$ from $y$ to $x$, retain its
specified reference $T$. On a fixed small conjugation-invariant
neighborhood of the identity for every $U_jT^{-1}$, put
\[
 P(U_1,\ldots,U_n)
   =\Phi(U_1T^{-1},\ldots,U_nT^{-1})T,
\]
where $\Phi$ is the \emph{actual selected local source mean}
on the reference-normalized tuple, with its original contour
multiplicities, ordering and reference. Right endpoint covariance
proves this representation on the original domain, and simultaneous
conjugation covariance follows from the same source covariance.
On the complementary set put $P=T$. Thus the completion is
measurable, group valued and endpoint covariant everywhere,
and equals the chosen analytic mean near the identity. No
different mean or additional mean equation is substituted.

Similarly, for each primitive ordinary average, use its original
logarithmic formula when all its reference-normalized transported
strings lie in the chosen small logarithm neighborhood, and its
reference string otherwise. The condition is conjugation invariant.
This gives a measurable gauge-covariant completion of
$\mathcal M_\ell$, hence of its composition $\mathcal F_R$.
The values outside the analytic domain are part of this proposed
full blocking map; equality with a different large-field completion
is not claimed. The following local formulas agree wherever
the original analytic maps agree.

Let $\mathcal H_R$ be the group of current-site gauges equal to
one at final block centers, with normalized independent Haar
measure at the $n_s$ other sites. For $g>0$ define
\begin{equation}\label{cmi-gauge-constant}
 c(g)=\int_G e^{-v(h)/g^2}\,dh>0.
\end{equation}
Then
\begin{equation}\label{cmi-orbit-identity}
 \int_{\mathcal H_R}e^{-\mathcal G_R(V^u)/g^2}\,du
       =c(g)^{n_s}.
\end{equation}
Indeed, by \eqref{cmi-covariance}, the factor at $(y,x)$ is
$e^{-v(P_R(V;y,x)u(x)^{-1})/g^2}$, whose integral is $c(g)$
by Haar invariance, independently of $V$. The site integrals
factor. There is no omitted gauge determinant: this is direct
integration over compact Haar variables with the displayed
normalization.

For an integrable gauge-invariant current density $\rho$, and
any bounded test function $f$ of the final configuration,
\begin{equation}\label{cmi-pushforward}
 \int dV\,\rho(V)f(\mathcal F_R(V))
 =c(g)^{-n_s}\int dV\,\rho(V)f(\mathcal F_R(V))
                            e^{-\mathcal G_R(V)/g^2}.
\end{equation}
Insert \eqref{cmi-orbit-identity}, use Fubini, then change
$V$ by the fixed gauge $u$; link Haar, $\rho$ and $\mathcal F_R$
are unchanged. This is an identity of finite pushforward
measures, requiring no globally regular delta density.
It also holds for a finite gauge-invariant current measure,
and for a gauge-invariant orbit-domain restriction.

A noninvariant small-field window is introduced \emph{after}
this insertion. Splitting the gauged integral into $\chi_{\rm sf}$
and $1-\chi_{\rm sf}$ is exact for any measurable
$0\le\chi_{\rm sf}\le1$ on the fiber. If $\rho$ is positive
both pieces are positive on real fields. Neither piece is
asserted to satisfy \eqref{cmi-pushforward} without the gauge
factor. The complement requires a separate large-field estimate.

\subsection{Identification of the physical centered branch}
\label{cmi-branch-section}

Let $U_f(W)$ be the proved full ordinary minimizer at total
refinement $MR$, with final spacing $R$ in current units, and set
\begin{equation}\label{cmi-pregauge}
 V_p(W)=\mathcal M_\ell^k(U_f(W)),\qquad
 \mathcal F_R(V_p(W))=W.
\end{equation}
Let $U_c(V)$ be the preceding-depth ordinary minimizer and
$S_M(V)=S(U_c(V))$. Choose a common current comparison radius
$\delta_R$ and then a final datum radius $b_R$ so that
$U_f(W)$ belongs to the current minimum class, every $U_c(V)$
with $\|\log V/i\|_\infty<\delta_R$ belongs to the full
final-spacing-$R$ minimum class, and the partial average and
gauges below stay within all common chart and invariance domains.
These are finitely many strict smallness conditions with
$R$-dependent constants, uniform in $M$ and volume. In dimension
four the coupling-free Wilson action has no additional scale
factor; the potential norms acquire only fixed powers of $R$.

The two actual comparisons are
\begin{equation}\label{cmi-two-comparisons}
 S_M(V_p)\le S(U_f)
       \le S(U_c(V_p))=S_M(V_p).
\end{equation}
For the second inequality exact composition makes $U_c(V_p)$
a full feasible competitor. For any other current $V$ in the
comparison ball with $\mathcal F_R(V)=W$, the same argument gives
\begin{equation}\label{cmi-current-minimum}
 S_M(V)\ge S(U_f)=S_M(V_p).
\end{equation}
This is a common-domain minimum comparison, not a
dynamic-programming assertion for arbitrary backgrounds.

Define the root-fixed current gauge and background by
\begin{equation}\label{cmi-rooted-background}
 u(y)=1,\qquad u(x)=P_R(V_p;y,x),\qquad V_*(W)=V_p(W)^u.
\end{equation}
Then $\mathcal F_R(V_*)=W$ and $P_R(V_*;y,x)=1$ exactly.
At fixed $R$ this is local analytic with
$\|\log u/i\|_\infty\le C_R\|\log V_p/i\|_\infty$.
The already proved coarse-action gauge invariance, with its
smooth fine point-value extension, implies $S_M(V_*)=S_M(V_p)$
after shrinking the common radii. Since $\mathcal G_R\ge0$
and $\mathcal G_R(V_*)=0$ on real fields, $V_*$ is a local
constrained minimum of $S_M+\mathcal G_R$.

The preceding Gaussian module supplies its uniformly positive
flat retained precision. The local analytic chart and
$D^2S_M=q_M+O(\|\log V\|_\infty)$ show that the pulled-back
Hessian on a common convex retained ball is
$A_{M,0}+O_R(\|w\|_\infty+\|z\|_\infty)$. It remains positive
after a fixed shrink. Therefore the centered real critical
point is unique in that common neighborhood; gauge zero
alone would not prove uniqueness. Constraint, gauge-zero and
first-derivative identities continue holomorphically to complex
$W$. No complex minimum is asserted.

The supported pre-gauge theorem gives the stronger amplitude
bound
\begin{equation}\label{cmi-pregauge-amplitude}
 a_p=\log V_p/i,\qquad
 \|a_p\|_\infty\le C_a\|w\|_\infty/R.
\end{equation}
Its supported-family constant is fixed at primitive $\ell$
independently of $s,k$. There is no corresponding $1/R$
contraction claim for the rooted potential $\log V_*/i$.

\subsection{The local density, Haar factors and pivot determinants}
\label{cmi-density-section}

Use Hermitian Lie coordinates orthonormal for $\tau$, of
dimension $q=\dim\mathfrak g$. Let $n_e$ and $n_c$ count
positive current and final macro bonds, and put
$n_z=q(n_e-n_c)$. In exponential coordinates write normalized
Haar as
\begin{equation}\label{cmi-Haar}
 d(e^{iX})=h_0j(X)\,dX,\qquad h_0>0,\quad j(0)=1.
\end{equation}
Consequently
$\delta_G(e^{iF})=h_0^{-1}\delta_{\mathbb R^q}(F)$,
with the coordinate measure understood: $j(F)$ is evaluated
at $F=0$.

Fix a real small $W$ and $V=V_*(W)$. The composite relative
constraint and its pivot derivative are
\begin{equation}\label{cmi-relative}
 F_V(B)=\log[\mathcal F_R(e^{iB}V)\mathcal F_R(V)^{-1}]/i,
 \qquad A_V=D F_V(0),\qquad
 A_{c,p(c)}=R^{1-d}\pi_c(V).
\end{equation}
Let $\psi_V$ be the ambient straightening, $K_V$ the linear
retained insertion, and $\Phi_V(z)=\psi_V(K_Vz)$.
Each retained component of $\Phi_V$ equals $z$ exactly.
Order ambient variables as retained coordinates followed by
pivots. The derivative of $\psi_V$ is block triangular;
its pivot diagonal blocks satisfy
\[
 D_{p(c)}F_V(\psi_V(B);c)\,
 D_{B(p(c))}\psi_V(B)(p(c))=A_{c,p(c)}.
\]
Thus
\begin{equation}\label{cmi-ambient-Jacobian}
 \det D\psi_V(B)=
   \prod_c\frac{\det A_{c,p(c)}}
                  {\det D_{p(c)}F_V(\psi_V(B);c)}.
\end{equation}
Let $j_c^{\rm str}(V;B)$ be the analytic logarithm of the
$c$th ratio, zero at $B=0$, and put
\begin{equation}\label{cmi-background-delta}
 \mathcal D_R(V)=\prod_c(\det A_{c,p(c)})^{-1}.
\end{equation}
The determinants have their positive flat orientation on the
real small chart.

Choose the small-field sector to be the component of this
relative chart with retained coordinates in a specified product
window, retaining any additional indicator explicitly. Changing
the square ambient variables and then solving $A_VB=0$ gives
its exact density
\begin{align}
 &c(g)^{-n_s}h_0^{n_e-n_c}\mathcal D_R(V)
 \int dz\,
    \prod_bj(\Phi_V(z)(b))\,
    e^{\sum_cj_c^{\rm str}(V;K_Vz)}
    \rho(e^{i\Phi_V(z)}V)
    e^{-\mathcal G_R(e^{i\Phi_V(z)}V)/g^2}.
                                                        \label{cmi-local-density}
\end{align}
The integration domain and every extra indicator are those
of the chosen sector. This uses the square Jacobian
\eqref{cmi-ambient-Jacobian}, not a determinant of the
rectangular derivative $D\Phi_V$. Equivalently the unrearranged
pivot density is the product of inverse
$D_{p(c)}F_V(\Phi_V(z);c)$ determinants.

Here the globally gauge-invariant $\rho$ in
\eqref{cmi-pushforward} is an input density or measure.
Only on the common small chart do we use a representation
\[
 \rho(U)=\exp[-S_M(U)/g^2+E(U)],
\]
with any field-independent factor tracked separately.
The local classical branch $S_M$ has not been promoted
to a global action on all configurations.

Put $z=g\xi$ for real $g>0$, and let $A=A_M(V)$ be the exact
retained Hessian of $S_M+\mathcal G_R$. Define
\begin{align}
 \mathcal R_3(V,z)
   &=S_M(e^{i\Phi_V(z)}V)+\mathcal G_R(e^{i\Phi_V(z)}V)
            -S_M(V)-\tfrac12\langle z,Az\rangle,\notag\\
 I(V,g,\xi)
   &=-g^{-2}\mathcal R_3(V,g\xi)
      +\sum_cj_c^{\rm str}(V;K_Vg\xi)
      +\sum_b\log j(\Phi_V(g\xi)(b))\notag\\
   &\hspace{1cm}+E(e^{i\Phi_V(g\xi)}V)-E(V).
                                                        \label{cmi-interaction}
\end{align}
There is no linear term because the background is the physical
centered branch. Formula \eqref{cmi-local-density} becomes
\begin{align}
 &c(g)^{-n_s}h_0^{n_e-n_c}\mathcal D_R(V)\,
   g^{n_z}(2\pi)^{n_z/2}\det(A)^{-1/2}
   e^{-S_M(V)/g^2+E(V)}\notag\\
 &\hspace{3cm}\times
   \mathbb E_{A^{-1}}\!
      \left[\chi_{\rm window}(\xi)e^{I(V,g,\xi)}\right].
                                                        \label{cmi-Gaussian-integral}
\end{align}
The Gaussian convention is the retained Euclidean density
$\det(A)^{1/2}(2\pi)^{-n_z/2}
e^{-\langle\xi,A\xi\rangle/2}$.
The covariance is that of this exact classical
action-plus-gauge Hessian; Haar remains in
\eqref{cmi-interaction}. The factors $c(g)^{-n_s}$,
$h_0^{n_e-n_c}$ and $g^{n_z}$ are vacuum normalizations.
$\mathcal D_R(V)$ and $\det A$ generally depend on the
background and are not vacuum constants.

\subsection{The literal retained cutoff}
\label{cmi-window-section}

For the original retained link/color balls set
\[
 \chi_{\rm window}(\xi)=
 \prod_{b\ {\rm retained}}\mathbf1_{\{|\xi(b)|\le T\}}.
\]
Since $\Phi_V(g\xi)(b)=g\xi(b)$ at every retained bond,
this is exactly the relative-coordinate window of the
chosen sector. It is not a cutoff on $K_V\xi$, or a
Euclidean ball on a whole macro fiber. Pivot, contour,
plaquette or reconstructed-field indicators, if included,
remain at their exact $\Phi_V(g\xi)$ pullbacks.
The macro sector is not identified with an intersection
of the primitive windows of $s$ separate integrations.

At fixed $R$, the normalized pivot logarithms are local
$O_R(|g|T)$ terms. Inversion invariance gives
$\log j(X)=O(|X|^2)$, so Haar logarithms are local
$O_R(g^2T^2)$ terms. The gauge third-and-higher remainder
is local $O_R(|g|T^3)$. The internal contours read no pivot,
and therefore its full exact factor satisfies
\begin{equation}\label{cmi-gauge-retained}
 \mathcal G_R(e^{i\Phi_V(z)}V)
       =\mathcal G_R(e^{iE_{\rm free}z}V).
\end{equation}
On an independent background product ball its constant
and linear terms must also be separated away from the
physical gauge-zero branch. All these estimates have
finite local supports and finite anchor multiplicities
at fixed $R$, without a sum over the old refinement $M$.
The classical action and inherited interaction require
their own support estimates.

\subsection{Pre-gauge continuation of the full inherited insertion}
\label{cmi-continuation-section}

With $V_*=V_p^u$, set, link by link,
\begin{equation}\label{cmi-prefluctuation}
 B_p(w,\xi)(b)=
   \operatorname{Ad}_{u(w;b_-)^{-1}}
                  \Phi_{V_*(w)}(g\xi)(b).
\end{equation}
Direct multiplication proves
\begin{equation}\label{cmi-fluctuation-identity}
 \bigl(e^{i\Phi_{V_*}(g\xi)}V_*\bigr)^{u^{-1}}
             =e^{iB_p(w,\xi)}V_p(w).
\end{equation}
The coupling remains inside the nonlinear chart:
$\Phi_V(g\xi)$ is not replaced by $g\Phi_V(\xi)$.
The gauge depends on $w$, not on $\xi$. Its real adjoint
is a color isometry; on the small complex gauge ball it
has a fixed $R$-dependent bound.

Let $E$ be a holomorphic total inherited functional on a
current potential product ball of radius $b$, with its
local gauge invariance on the original common domain.
Individual activity coefficients need not be gauge invariant.
On a sufficiently small common overlap, both arguments of
\eqref{cmi-fluctuation-identity}, and a small connecting
gauge path, are in this domain. Their $E$ values therefore
agree, first on real fields and then by holomorphic uniqueness.
If only $e^E$ is supplied, a compatible nonvanishing
holomorphic logarithm branch must be supplied before making
the corresponding statement about $E$.

On a fixed target $b$-ball below all classical domains,
\eqref{cmi-pregauge-amplitude} and the local chart give
\begin{equation}\label{cmi-common-ball}
 \left\|\frac1i\log[e^{iB_p(w,\xi)}V_p(w)]\right\|_\infty
       \le C_a b/R+C_R|g|T.
\end{equation}
Choose $R\ge4C_a$ and $C_R|g|T\le b/4$.
The pre-gauge argument then stays inside the old $b$-ball
uniformly in $M$ and volume. Its inherited functional is
defined and holomorphic on the whole target ball, even
where the rooted argument falls outside the original
domain of $E$. It is a continuation of the common-domain
germ, not an assertion that $E$ was already defined at
that rooted argument. Replacing $g\xi$ by $\lambda g\xi$
gives a complex $\lambda$-disc of radius at least
$cb/(C_R|g|T)$ within the strict chart margins.

The \emph{entire} inherited contribution in
\eqref{cmi-Gaussian-integral} is continued: replace its
outside $E(V_*)$ by $E(V_p)$, and replace the inside
increment by
\[
 E(e^{iB_p(w,\xi)}V_p(w))-E(V_p(w)).
\]
Leaving $E(V_*)$ outside would be undefined on part of
the enlarged domain. There is no change of integration
variable $\xi$. The Gaussian, its determinant, the pivot
and Haar factors, the absolute gauge penalty and the
retained cutoff remain exactly those of
\eqref{cmi-Gaussian-integral}; no adjoint Jacobian or
transformed window is introduced.

\subsection{What holomorphy does and does not imply}
\label{cmi-scope-section}

At fixed finite regulator and fixed real retained window,
the continued integrand is holomorphic in $w$. The proved
complex precision neighborhood supplies a consistent
analytic determinant square root. Every compact subdomain
has an integrable uniform majorant, so the fixed-window
integral is holomorphic by differentiation under the
integral. Equality on the nonempty overlap identifies it
as the continuation of the original small-field germ.

For real data and positive physical input, the exact
integrand is positive and the retained window has positive
measure. The integral is then strictly positive. This
does not imply nonvanishing on a volume-independent
complex ball: an integral of pointwise nonzero analytic
factors can vanish. A finite-volume logarithm germ is
available near a positive point; a uniform complex
logarithm needs the full cluster hypotheses and budgets.

Similarly, pointwise $\lambda$-Cauchy applied to a total
$E$ can give only an extensive estimate. An anchored
interaction bound needs a strict joint $(w,\xi)$ family
for the pre-gauge argument, with remote background
supports and an explicit local fluctuation collar.
That separate support module is not inferred from
\eqref{cmi-common-ball}. Purely classical coefficients
have independent fixed classical domains; the quantum
radius may be chosen below them once, rather than
reused as their outer analytic radius at every step.

The present module proves the finite-regulator
normalization and local density identities, identifies
the physical centered branch, preserves the exact
window, and supplies the pointwise pre-gauge extension.
It does not establish a stable quantum iteration,
uniform nonvanishing without the cluster bounds,
large-field control, mixed-history matching, coupling
flow or a continuum measure.

\section{Place in the construction}

The longer physical blocking step is now defined, including the
source chart, its centered Gaussian, its gauge insertion and its
small-field cutoff. The pre-gauge representative of Notebook 25
belongs to this actual composite constraint. Choosing $R$ after its
primitive-base constants have been fixed can therefore provide a
strict amplitude margin for an inherited field on a fixed analytic
ball. This is the role of the new macrostep.

The next analytic task is to control the full inherited insertion
as strict joint activities in the coarse background and retained
fluctuations. A local-core expansion can keep its nearby field
dependence unexpanded and spend Taylor estimates only on its distant
tail. The precision and purely classical interactions can be
expanded directly in the coarse background before Gaussian
integration, using their existing finite classical support reserves.
Both operations require their explicit norm budgets; the pointwise
continuation proved here is not a substitute.

Even a common analytic domain and a convergent one-step logarithm
would leave a noncontracting activity-amplitude recurrence.
Extracting and controlling the vacuum and marginal gauge-action
terms is required for an arbitrarily long quantum iteration.
Large-field complements, mixed histories, observable estimates and
continuum reconstruction remain parts of the existence argument.

\begin{thebibliography}{9}
\bibitem{balaban109cm}
T.~Ba\l aban,
\emph{Renormalization group approach to lattice gauge field
theories. I},
Commun. Math. Phys. \textbf{109} (1987), 249--301.
\href{https://scispace.com/pdf/renormalization-group-approach-to-lattice-gauge-field-3ygai0vxt1.pdf}{Original article copy}.
The local ordinary average and contour mean are taken from
pp.~252--254; the composite-step arguments are developed here.
\bibitem{priorcm}
\emph{Toward a four-dimensional Yang--Mills construction},
research Notebooks 05--26 in this workspace.
The classical minimum and gauge identities, flat quadratic
locality, strict classical support families and the supported
pre-gauge representative are the previously established inputs.
\bibitem{sourcecm}
E.~Sussman and collaborators,
\href{https://github.com/ethanwsussman/Yang--Mills-LLM-experiment}{Yang--Mills-LLM-experiment},
deposited raw accounts at commit
\texttt{e392ef45013a581664bab76b90867170b30fe10e}.
\end{thebibliography}
\end{document}
